% ============================================================
% Template LaTeX ABNT NBR 14724/2024 — IFB
% Briefing de Projeto — Luminária de Mesa
% Atividade 2 — Metodologia e Prática de Projeto
% ============================================================

\documentclass[
  12pt,
  a4paper,
  oneside,
  brazil
]{article}

% ── Codificação e Idioma ──────────────────────────────────────────────────────
\usepackage[T1]{fontenc}
\usepackage[utf8]{inputenc}
\usepackage[brazil]{babel}

% ── Fonte Times New Roman ─────────────────────────────────────────────────────
\usepackage{mathptmx}

% ── Geometria da página — margens ABNT anverso ────────────────────────────────
\usepackage[
  a4paper,
  left=3cm,
  top=3cm,
  right=2cm,
  bottom=2cm,
  headheight=14pt,
  headsep=0.5cm
]{geometry}

% ── Espaçamento 1,5 entre linhas (ABNT) ──────────────────────────────────────
\usepackage{setspace}
\onehalfspacing

% ── Recuo do parágrafo: 1,25cm (ABNT) ────────────────────────────────────────
\usepackage{indentfirst}
\setlength{\parindent}{1.25cm}
\setlength{\parskip}{0pt}

% ── Formatação de seções — numeração progressiva NBR 6024 ────────────────────
\usepackage{titlesec}

\titleformat{\section}
  {\normalfont\normalsize\bfseries\MakeUppercase}
  {\thesection\quad}
  {0em}{}

\titleformat{\subsection}
  {\normalfont\normalsize\bfseries}
  {\thesubsection\quad}
  {0em}{}

\titleformat{\subsubsection}
  {\normalfont\normalsize\bfseries\itshape}
  {\thesubsubsection\quad}
  {0em}{}

\titlespacing*{\section}      {0pt}{18pt}{12pt}
\titlespacing*{\subsection}   {0pt}{18pt}{6pt}
\titlespacing*{\subsubsection}{0pt}{12pt}{6pt}

% Seções primárias iniciam em nova página
\let\oldsection\section
\renewcommand{\section}{\clearpage\oldsection}

% ── Numeração de páginas — canto superior direito, 10pt ──────────────────────
\usepackage{fancyhdr}
\pagestyle{fancy}
\fancyhf{}
\fancyhead[R]{\fontsize{10}{12}\selectfont\thepage}
\renewcommand{\headrulewidth}{0pt}
\renewcommand{\footrulewidth}{0pt}

\fancypagestyle{plain}{
  \fancyhf{}
  \fancyhead[R]{\fontsize{10}{12}\selectfont\thepage}
  \renewcommand{\headrulewidth}{0pt}
}

% ── Listas ───────────────────────────────────────────────────────────────────
\usepackage{enumitem}
\setlist{nosep, leftmargin=*}
\providecommand{\tightlist}{\setlength{\itemsep}{0pt}\setlength{\parskip}{0pt}}

% ── Tabelas e Quadros ─────────────────────────────────────────────────────────
\usepackage{longtable}
\usepackage{booktabs}
\usepackage{array}
\renewcommand{\arraystretch}{1.3}

% ── Legendas ─────────────────────────────────────────────────────────────────
\usepackage{caption}
\DeclareCaptionLabelSeparator{emdash}{\,---\,}
\DeclareCaptionFormat{abnt}{\fontsize{10}{12}\selectfont#1#2#3}
\captionsetup{
  format=abnt,
  justification=centering,
  singlelinecheck=false,
  labelsep=emdash,
  position=above,
  skip=3pt
}

% ── Sumário ───────────────────────────────────────────────────────────────────
\usepackage{tocloft}
\renewcommand{\contentsname}{SUMÁRIO}

% ── Quadros ───────────────────────────────────────────────────────────────────
\newcommand{\quadrotitle}[2]{%
  \vspace{1em}%
  \noindent\textbf{Quadro #1 --- #2}\par\vspace{0.3em}%
}
\newcommand{\quadrofonte}[1]{\par\noindent\textit{Fonte: #1}\vspace{1em}\par}

% ── URL e Hyperlinks ─────────────────────────────────────────────────────────
\usepackage{url}
\PassOptionsToPackage{hyphens}{url}
\usepackage[hidelinks, colorlinks=false, breaklinks=true]{hyperref}

% ── Figuras ──────────────────────────────────────────────────────────────────
\usepackage{graphicx}
\usepackage{float}

\date{}

% ============================================================
\begin{document}
% ============================================================

% ══════════════════════════════════════════════════════════════════════════════
% CAPA — NBR 14724/2024
% ══════════════════════════════════════════════════════════════════════════════
\begin{titlepage}
  \thispagestyle{empty}
  \begin{center}

    \includegraphics[width=0.25\textwidth]{../../../logo_ifb_despro.png}\\[1.8em]

    {\normalsize\bfseries
      Instituto Federal de Educação, Ciência e Tecnologia de Brasília --- IFB\\[0.3em]
      \textit{Campus} Samambaia\\[0.3em]
      Curso Superior de Tecnologia em Design de Produto
    }

    \vfill

    {\large\bfseries
      BRIEFING DE PROJETO --- LUMINÁRIA DE MESA
    }\\[1.2em]

    {\normalsize
      Atividade 2: Briefing e Painel Semântico
    }

    \vfill

    {\normalsize Brasília/DF\\2026}

  \end{center}
\end{titlepage}
\setcounter{page}{2}

% ══════════════════════════════════════════════════════════════════════════════
% FOLHA DE ROSTO — NBR 14724/2024
% ══════════════════════════════════════════════════════════════════════════════
\begin{titlepage}
  \thispagestyle{empty}
  \begin{center}

    {\normalsize\bfseries VICTOR COSTA · ANNA JI · LUÍS COSTA}

    \vfill

    {\large\bfseries BRIEFING DE PROJETO --- LUMINÁRIA DE MESA}

    \vfill

    \begin{flushright}
      \begin{minipage}{0.5\textwidth}
        \normalsize
        Atividade apresentada ao Curso Superior de Tecnologia em Design de Produto
        do Instituto Federal de Educação, Ciência e Tecnologia de Brasília --
        Campus Samambaia, como requisito parcial para a aprovação na disciplina
        de Metodologia e Prática de Projeto.

        \vspace{1em}
        Professora: Profª. Fernanda Freitas Costa de Torres
      \end{minipage}
    \end{flushright}

    \vfill

    {\normalsize Brasília/DF -- 2026}

  \end{center}
\end{titlepage}

% ══════════════════════════════════════════════════════════════════════════════
% RESUMO — NBR 6028/2003
% ══════════════════════════════════════════════════════════════════════════════
\clearpage
\thispagestyle{empty}
\begin{center}
  {\normalsize\bfseries RESUMO}
\end{center}

\vspace{1em}
\noindent
Este documento apresenta o briefing de projeto para o desenvolvimento de uma
luminária de mesa concebida como presente de design para mulher de classe
média-alta/alta, com aproximadamente 40 anos, moradora do bairro Noroeste, em
Brasília. O briefing define o problema de projeto, o público-alvo, os contextos de
uso e os requisitos que orientarão as etapas de pesquisa, conceituação e
desenvolvimento do produto.

\vspace{1em}
\noindent\textbf{Palavras-chave:} luminária de mesa; briefing; design de produto;
presente de design; Noroeste Brasília.

% ══════════════════════════════════════════════════════════════════════════════
% SUMÁRIO — NBR 14724/2024
% ══════════════════════════════════════════════════════════════════════════════
\clearpage
\thispagestyle{empty}
\tableofcontents

% ══════════════════════════════════════════════════════════════════════════════
% CONTEÚDO
% ══════════════════════════════════════════════════════════════════════════════

\section{Identificação do Projeto}

\noindent\textbf{Projeto:} Luminária de Mesa --- Presente de Design\\[0.3em]
\noindent\textbf{Equipe:} Victor Costa, Anna Ji, Luís Costa\\[0.3em]
\noindent\textbf{Disciplina:} Metodologia e Prática de Projeto\\[0.3em]
\noindent\textbf{Professora:} Profª. Fernanda Freitas Costa de Torres\\[0.3em]
\noindent\textbf{Curso:} Tecnologia em Design de Produto --- IFB Campus Samambaia\\[0.3em]
\noindent\textbf{Data:} Abril de 2026

\section{O Problema}

Desenvolver uma luminária de mesa que funcione como presente de design para mulher
de classe média-alta/alta, com aproximadamente 40 anos, moradora do bairro Noroeste
em Brasília. O produto precisa operar em dois registros ao mesmo tempo: ser
reconhecível como presente especial no momento da entrega e justificar sua presença
no espaço doméstico no uso cotidiano.

O desafio está em conciliar valor percebido, funcionalidade e viabilidade de produção
sem que nenhum desses três critérios comprometa os outros. Um objeto com alto apelo
visual mas luz ruim falha no uso; um objeto tecnicamente impecável mas genérico
falha como presente.

\section{Contexto}

O mercado de presentes de design opera com uma mediação dupla: quem compra e quem
usa são pessoas diferentes, com critérios de avaliação diferentes. O comprador decide
no ponto de venda ou catálogo, com base em aparência, embalagem e reputação da peça.
O usuário decide ao longo do tempo, com base em desempenho, conforto e adequação
ao espaço.

O Noroeste é um bairro planejado de Brasília, inaugurado a partir de 2012, com perfil
socioeconômico elevado e arquitetura predominantemente contemporânea. Os apartamentos
da região têm plantas abertas, pé-direito generoso e acabamentos de alto padrão. O
habitante típico tem entre 30 e 55 anos, formação universitária, renda elevada e
repertório estético construído ao longo de uma trajetória de consumo criterioso.
Pesquisa antes de comprar e espera coerência entre o objeto e o ambiente em que ele
vai habitar.

A mulher de aproximadamente 40 anos que configura o público-alvo já passou pela
fase das aquisições utilitárias. O apartamento está mobilado; o que ela busca agora
são peças com identidade, que acrescentem camadas ao espaço. Uma luminária, nesse
contexto, compete com outros objetos de curadoria doméstica por presença e atenção.

\section{Público-Alvo}

\noindent\textbf{Perfil demográfico:} Mulheres entre 35 e 50 anos, classes A e B,
moradoras de apartamentos contemporâneos em Brasília, com preferência pelo Noroeste
e regiões de perfil equivalente.

\vspace{0.5em}
\noindent\textbf{Perfil profissional e econômico:} Profissionalmente ativas, com
formação superior e renda mensal acima de R\$ 10.000. Tomam decisões de compra com
autonomia e são o público primário ou o destinatário direto do produto.

\vspace{0.5em}
\noindent\textbf{Relação com objetos:} Valorizam autoria, qualidade de execução e
coerência estética. Não são atraídas por produtos de apelo massivo nem por tendências
de ciclo curto. Têm disposição para pagar mais quando o objeto justifica o preço.

\vspace{0.5em}
\noindent\textbf{Espectro de compradores no contexto de presente:} Parceiros, filhos
adultos, amigos próximos e colegas com nível socioeconômico equivalente, em busca de
um presente que comunique consideração e gosto sem recorrer ao óbvio.

\section{Contexto de Uso}

A luminária será usada em ambiente residencial, predominantemente em superfícies de
apoio em três cômodos:

\begin{itemize}
\tightlist
\item \textbf{Quarto:} sobre criado-mudo ou mesa lateral, como iluminação de leitura
  e presença noturna;
\item \textbf{Home office:} sobre escrivaninha, como iluminação de trabalho com
  conforto visual;
\item \textbf{Sala de estar:} sobre aparador, console ou mesa de canto, como peça de
  composição e iluminação de acento.
\end{itemize}

Em todos os cenários, a luminária divide superfície com outros objetos: livros,
plantas, objetos de coleção. O produto precisa funcionar bem nessa composição sem
subordinar tudo ao seu entorno nem dominar a cena.

A situação de presente acrescenta um momento de entrega: o produto precisa comunicar
valor antes mesmo de ser usado, o que tem implicações para embalagem e apresentação,
ainda que fora do escopo desta fase.

\section{Requisitos Funcionais}

\begin{itemize}
\tightlist
\item Iluminação de mesa para uso residencial, adequada à leitura e ao estar;
\item Temperatura de cor entre 2.700 K e 3.000 K (branco quente);
\item Controle de intensidade (dimerização) desejável;
\item Cabo com comprimento mínimo de 1,8 m, com plug compatível com padrão brasileiro
  (NBR 14136);
\item Base estável em superfície plana, sem necessidade de fixação;
\item Peso compatível com transporte manual e reposicionamento frequente.
\end{itemize}

\section{Requisitos de Negócio}

\begin{itemize}
\tightlist
\item Posicionamento de preço de venda ao público entre R\$ 900 e R\$ 1.800;
\item Produção em pequenos lotes ou sob encomenda; não é objetivo desta fase projetar
  para linha de produção em larga escala;
\item Produto deve ser viável para produção por fornecedores nacionais.
\end{itemize}

\section{Restrições}

\begin{itemize}
\tightlist
\item Não utilizar fontes de iluminação fluorescentes ou incandescentes;
\item Componentes elétricos devem atender à ABNT NBR 5410 e ser certificados pelo
  Inmetro;
\item Não incluir funcionalidades de conectividade ou automação nesta fase do projeto;
\item Soluções que dependam de importação de componentes ou processos sem viabilidade
  de replicação no Brasil estão fora do escopo.
\end{itemize}

\section{Referências}\label{referencias}

BÜRDEK, Bernhard E. \textbf{Design}: história, teoria e prática do design de produtos.
São Paulo: Edgard Blücher, 2006.

MORAES, Dijon de. \textbf{Metaprojeto}: o design do design. São Paulo: Edgard Blücher,
2010.

NORMAN, Donald A. \textbf{Design emocional}: por que adoramos (ou detestamos) os objetos
do dia a dia. Rio de Janeiro: Rocco, 2008.

\end{document}
